\section{综合训练与复盘}
\subsection{练习方式}
本章三道题均为教学训练题，不对应某次认证的原版试卷。建议先不看解答，给自己安排30、50、90分钟的独立尝试，再用解答检查建模过程。用时只是训练建议，不是官方考试分配。

\example{训练A 差值数对}
给 n 个整数和非负整数 c，求满足 $i<j$ 且 $|a_i-a_j|=c$ 的位置对数量。$1\le n\le2\times10^5$，$|a_i|\le10^9$，$0\le c\le2\times10^9$。第一行 n、c，第二行数组，输出一个整数。
\begin{verbatim}
输入
5 2
1 3 1 3 3
输出
6
\end{verbatim}
\textbf{解答。}先统计每个值的频次。c>0时，只从较小值x寻找x+c，不会重复统计同一无序数对；贡献是两组频次乘积。c=0时，频次k的组中任选两个位置，贡献 $k(k-1)/2$。答案最多为 $n(n-1)/2$，需要long long。时间 $O(n\log n)$，空间 $O(n)$。
\Code{24_pairs.cpp}
\textbf{检查用例。}数组 $[7,7,7]$、c=0应输出3；单元素时应输出0。若输出9，说明把相同位置配对或对称重复都算了进去。

\example{训练B 矩形增量的最大值}
初始有 $n\times m$ 全0矩阵，执行q次操作，每次给闭矩形左上 $(x_1,y_1)$ 到右下 $(x_2,y_2)$ 加v。最终输出最大元素值，以及等于最大值的格子数。$1\le n,m\le1000$，$0\le q\le2\times10^5$，$|v|\le10^6$，坐标合法。第一行n、m、q，随后每行五个数。
\begin{verbatim}
输入
3 3 2
1 1 2 2 2
2 2 3 3 3
输出
5 1
\end{verbatim}
\textbf{解答。}把每次矩形更新记录在二维差分四角，最后恢复每格真实值，同时维护最大值和次数。只统计1到n行、1到m列，不能把辅助边界算进去。允许负更新，最大值要从负无穷开始，不能默认0。时间 $O(nm+q)$，空间 $O(nm)$。
\Code{25_rectangle.cpp}
\textbf{检查用例。}q=0时输出 \texttt{0 n*m}，第二项应计算成实际格子数；若整个矩阵一次减5，则最大值为-5而不是0。

\example{训练C 非抢占任务调度}
有一台处理器和n个任务，每个任务给定到达时刻、优先级和执行时间，输入顺序就是任务编号。处理器空闲时，在所有已经到达但尚未执行的任务中选择优先级最高者；并列时选到达更早者，再并列选编号更小者。一旦开始执行就\textbf{不能被抢占}。

一个任务完成时，必须先把恰好在该时刻到达的任务加入候选，再选择下一项。若没有已到达任务，时钟直接跳到下一次到达。输入不保证按到达时刻有序。$1\le n\le2\times10^5$，到达时刻和优先级在0到 $10^9$ 之间，执行时间在1到 $10^9$ 之间。

第一行n，随后每行依次给到达时刻、优先级、执行时间。按输入编号输出每个任务的开始和结束时刻，最后一行输出最大等待时间，等待时间为开始减到达。
\begin{verbatim}
输入
4
0 1 5
1 1 2
2 3 1
5 2 2
输出
0 5
8 10
5 6
6 8
7
\end{verbatim}
\textbf{推演。}任务1在0到5执行，中途到达的高优先级任务不能打断它。时刻5，先加入任务4；待选任务2、3、4中任务3优先级最高，所以执行顺序为1、3、4、2。任务2等到8开始，等待7。

\textbf{状态划分。}排序数组保存还未到达的任务，指针i指向下一项；优先队列保存已经到达的任务；now保存处理器可再次分配工作的时刻。每轮先补全所有到达时间不大于now的任务，再选堆顶。

\textbf{正确性。}选择前，堆恰好包含全部可选任务；比较器严格实现三级优先顺序。执行结束后更新now，不处理中途抢占，符合规则。每个任务入堆出堆各一次，排序与堆操作总时间 $O(n\log n)$，空间 $O(n)$。结束时刻可能达到 $10^{14}$ 量级，应使用long long。
\Code{26_scheduler.cpp}

\subsection{如何复盘一次错误}
先保存最小失败输入，写出期望输出，再归类原因。输入解析错误要检查分词和换行；状态错误要明确变量对应时刻；算法错误要重新验证单调性、转移或贪心理由；复杂度错误要估计最坏规模。修复后还要补一个能够稳定复现原错误的测试。

例如调度程序把新任务加入堆的条件写成“到达时刻小于now”，就会漏掉恰好在任务结束时到达的任务。用“任务1在0开始、执行2；任务2在2到达”即可定位边界。改成小于等于后，再检查已有等待任务与同时到达任务的优先顺序。

\subsection{提交前检查表}
\begin{enumerate}
\item 输出内容、顺序和格式是否与题意一致，有没有调试输出残留。
\item 最小输入、全相等、全负数、空结果、不可达等合法边界是否处理。
\item 数值和数组范围是否足够，r+1、i-1等辅助位置是否分配。
\item 每个循环能否终止；每个递归状态是否完整；是否可能爆栈。
\item 每个容器访问前是否检查为空或end；比较器是否严格。
\item 程序复杂度是否覆盖最大数据，而不仅是样例。
\end{enumerate}
\newpage
